\chapter{Probabilistic Foundations: Arrivals, Service, and Little's Law}\label{sec:L1}

This lecture defines the vocabulary of queueing theory and proves three tools that
every later lecture relies on: the memorylessness of the exponential distribution,
the Poisson process together with PASTA, and Little's law. We assume undergraduate
probability (conditional expectation, the strong law of large numbers, probability
generating functions).

\section{Queueing systems and performance measures}

\begin{definition}[Queueing system]\label{L1:def:system}
A \term{queueing system} consists of four elements.
\begin{enumerate}
\item \term{Customers} (or jobs): entities that arrive at the system, require some
  amount of processing, and leave. The arrival time of the $k$-th customer is $a_k$,
  numbered so that $0\le a_1\le a_2\le\cdots$.
\item The \term{arrival process}: the sequence of arrival times $(a_k)_{k\ge1}$, or
  equivalently the counting process $A(t)=\#\{k: a_k\le t\}$.
\item A \term{server} and \term{service times}: the server is the resource that
  processes customers. Customer $k$ needs time $S_k$ when it has the server to
  itself. $S_k$ is also called the \term{work} or service requirement of customer $k$.
\item A \term{waiting room} and a \term{service discipline}: the space where customers
  arriving at a busy server wait, and the rule that decides whom to serve, when, and
  how much.
\end{enumerate}
\end{definition}

\begin{definition}[Service disciplines]\label{L1:def:discipline}
\term{FIFO} (first in, first out) serves one customer at a time, to completion, in
order of arrival. \term{Processor sharing} (PS) serves all customers in the system
simultaneously at equal rates; it is defined formally in \cref{sec:L2}.
A discipline under which the server never idles while work is present is called
\term{work-conserving}. Both FIFO and PS are work-conserving.
\end{definition}

\begin{definition}[Kendall notation]\label{L1:def:kendall}
A queueing system is written $A/S/c/K/N$. Here $A$ is the distribution of the
interarrival times, $S$ the distribution of the service times, $c$ the number of
servers, $K$ the system capacity ($\infty$ if omitted), and $N$ the size of the
customer population ($\infty$ if omitted). The distribution symbols are $M$ (Markov:
exponential and independent), $D$ (deterministic: constant), and $G$ (general:
an arbitrary distribution, i.i.d.).
\end{definition}

For example, M/M/1 has Poisson arrivals, exponential service and one server
(\cref{sec:L2}); M/G/1 allows a general service distribution (\cref{sec:L3});
M/M/1//N is a closed system with only $N$ customers (\cref{sec:L4}); and
M/G/$\infty$ is an infinite-server system with no waiting (\cref{sec:L4}).

\begin{definition}[Performance measures]\label{L1:def:measures}
Let $d_k$ be the time at which customer $k$ leaves the system.
\begin{enumerate}
\item The \term{sojourn time} (response time) is $W_k=d_k-a_k$.
\item The \term{waiting time} $W_{q,k}$ is the time from the arrival of customer $k$
  until its service first begins. For a FIFO single server, $W_k=W_{q,k}+S_k$.
\item The \term{number in system} is $N(t)=\#\{k: a_k\le t<d_k\}$. The number in the
  waiting room is written $N_q(t)$.
\end{enumerate}
\end{definition}

Quantities in queueing theory are defined by two different kinds of averages, and the
distinction will be essential several times.

\begin{definition}[Time averages and customer averages]\label{L1:def:averages}
The \term{time average} of the process $N(t)$ is
$L=\lim_{t\to\infty}\frac1t\int_0^t N(s)\dd s$, and the \term{customer average} of the
sequence $(W_k)$ is $W=\lim_{n\to\infty}\frac1n\sum_{k=1}^n W_k$ (when the limits
exist). $L_q$ and $W_q$ are defined in the same way. If the limits are constants with
probability one, then in the stationary regime they equal $\E[N]$ and $\E[W]$
respectively.
\end{definition}

\begin{definition}[Arrival rate, load, stability]\label{def:load}
The \term{arrival rate} is $\lambda=\lim_{t\to\infty}A(t)/t$.
The \term{load} (utilisation) of a single-server system with mean service time
$\E[S]$ is
\[
  \rho=\lambda\,\E[S].
\]
$\rho$ is the amount of work arriving per unit time and is dimensionless.
The system is \term{stable} if $N(t)$ has a proper limiting distribution (or at least
$L<\infty$) and every customer leaves in finite time.
\end{definition}

A single server processes at most one unit of work per unit time, so if $\rho>1$ the
unfinished work grows linearly and the system cannot be stable. For a work-conserving
G/G/1 system (i.i.d.\ interarrival and service times) it is known that $\rho<1$ is
essentially necessary and sufficient for stability \cite{loynes1962}. For the models
in these notes we assume $\rho<1$ as the stability condition and verify it model by
model.

\section{From a serving system to a queueing model}\label{L1:sec:projection}

A queueing model of a serving system is a projection: it keeps what decides who waits
and for how long, and drops the rest. \Cref{fig:projection} shows the projection used
throughout these notes. The stations it names (processor sharing, delay stations,
closed loops) are defined formally in later lectures; here they are described only
informally, so that every later theorem has a place in the picture.

\begin{figure}[t]
\centering
\begin{tikzpicture}[x=0.86cm, >=Stealth, font=\small, line width=0.6pt,
  station/.style={circle, draw, thick, minimum size=12mm, align=center},
  note/.style={font=\footnotesize, align=center, text=black!75}]
  % memory admission (engine-dependent): FCFS waiting line and m slots
  \foreach \i in {0,1}
    \draw[fill=orange!10] (-4.1+0.55*\i,-0.35) rectangle ++(0.5,0.7);
  \draw (-4.15,-0.35) -- (-3.0,-0.35) (-4.15,0.35) -- (-3.0,0.35);
  \node[draw, thick, rounded corners=1pt, fill=orange!10, minimum width=15mm, minimum height=11mm] (S) at (-1.75,0) {};
  \foreach \x in {-0.45,0,0.45}
    \foreach \y in {-0.22,0.22}
      \draw[orange!60!black, fill=white] ($(S)+(\x-0.17,\y-0.15)$) rectangle ++(0.34,0.3);
  \node[note, anchor=south] (Slab) at (-2.55,0.75) {\textbf{Memory admission}\\$m$ slots, FCFS};
  % prefill queue: waiting slots and a FIFO server
  \foreach \i in {0,1,2}
    \draw[fill=blue!8] (0.55*\i,-0.35) rectangle ++(0.5,0.7);
  \draw (-0.05,-0.35) -- (1.65,-0.35) (-0.05,0.35) -- (1.65,0.35);
  \node[station, fill=blue!8] (P) at (2.75,0) {FIFO};
  \node[note, anchor=south] (Plab) at (1.9,0.75) {\textbf{Prefill} (Lecture 3)\\server = token budget};
  \node[note, anchor=north] (Pttft) at (2.75,-0.8) {TTFT $=W_q+S$};
  % decode batch as a PS station
  \node[station, minimum size=15mm, fill=red!8] (D) at (6.7,0) {};
  \foreach \a in {0,72,...,288}
    \fill[red!60!black] ($(D)+(\a:0.42)$) circle (1.5pt);
  \node[font=\footnotesize] at (D) {PS};
  \node[note, anchor=south] (Dlab) at (6.25,0.9) {\textbf{Decode} (Lecture 2)\\capacity $\varphi(m)$};
  % tool call as a delay station
  \node[draw, thick, rounded corners, fill=teal!8, minimum width=24mm, minimum height=11mm] (T) at (3.4,-2.8) {};
  \foreach \k in {-2,-1,0,1,2}
    \draw[teal!60!black] ($(T)+(0.42*\k,0)$) circle (0.16);
  \node[note, anchor=north] (Tlab) at ($(T.south)+(0,-0.08)$) {\textbf{Tool call}: delay station, no queue (Lecture 4)};
  % flows
  \draw[->] (-7.2,0) -- (-4.2,0);
  \node[note, anchor=south] at (-6.3,0.05) {new sessions\\Poisson $\Lambda$};
  \draw[->] (S.east) -- (-0.1,0);
  \draw[->] (P) -- node[note, above] {first token} (D);
  \draw[->] (D.east) -- (8.9,0) node[note, right] {session\\ends};
  \draw[->] (D.south) |- node[note, right, pos=0.3] {resume\\w.p.\ $p$} (T.east);
  \draw[->] (T.west) -| node[note, left, pos=0.8] (Nlab) {next turn\\(hit or miss)} (-3.6,-0.38);
  % the slot is held from admission to the end of decode
  \draw[->, dotted, thick, orange!60!black] (D.north east) -- ++(0.35,0.35) |- (-1.75,2.35)
    node[note, pos=0.75, above, text=orange!50!black] {slot freed when decode ends} -- (S.north);
  % memory pool
  \begin{scope}[on background layer]
    \node[draw, dashed, rounded corners=4mm, fill=black!3, inner sep=3mm,
          fit={(Slab) (Plab) (Dlab) (Tlab) (Nlab) (S) (P) (D) (T) (-4.15,0) (-1.75,2.6)}] (M) {};
  \end{scope}
  \node[note, anchor=south west, align=left] at ($(M.north west)+(0.1,0.05)$)
    {\textbf{Memory pool} (dashed): $M$ bytes: KV of running turns, and cached KV of paused sessions\\
     (evictable); a constraint in the paper's model, a queue of slots in engines that admit whole prompts};
\end{tikzpicture}
\caption{The serving system projected onto queues. Turns are the customers of the
stations; sessions are the customers of the closed loop through the tool call. The
dashed pool is shared memory: what it drops decides whether the next turn is a hit
(short prefill work) or a miss (long prefill work). The memory-admission station
(orange) is \emph{not} in the paper's model: it appears in engines that admit a turn
only when its whole context fits, in arrival order; a turn then holds a slot from the
start of its prefill to the end of its decode (\cref{L5:sec:honest}).}
\label{fig:projection}
\end{figure}

\paragraph{Customers at two levels.} A \emph{turn} is a customer of the prefill and
decode stations: it arrives, waits, is served and leaves. A \emph{session} is a customer
at a coarser level: it circulates through prefill, decode and the tool call, and leaves
only when the agent program ends. New sessions arrive from outside as a Poisson process
of rate $\Lambda$, but the turns of a live session do not arrive independently of the
state: the next turn exists only after the previous one has been served. The prefill and
decode lectures look at turns; the lectures on closed systems, memory and admission look
at sessions.

\paragraph{Three stations.} \emph{Prefill} is a single FIFO server whose capacity is the
token budget of an engine step: serving engines by default give that budget to the oldest
pending prefill first. With Poisson arrivals this is the M/G/1 queue of
\cref{sec:L3}; with a finite number $N$ of live sessions it is the closed system of
\cref{sec:L4}. \emph{Decode} is a processor-sharing station: every step advances every
turn in the batch by one token, so the batch shares the device equally, with a total
rate $\varphi(m)$ that depends on the batch size $m$ (\cref{sec:L2}). The \emph{tool call}
is a delay station: every session calls its own tool at once, without queueing, for a
random think time.

\paragraph{Memory is a constraint, and in some engines a queue.} The pool of $M$ bytes
holds the KV cache of running turns and of paused sessions. In the paper's model
nobody queues ``at'' the memory: the pool only bounds how much state can stay resident,
and counting how much is resident on average is a calculation about sessions that stay
in the system for a random time, which \cref{sec:L4} does with the M/G/$\infty$ model
and Campbell's theorem. Many engines, however, admit a waiting turn only when its whole
context fits, and in arrival order; vLLM does so by default (its scheduler checks that
the full input sequence fits before admitting a request, and stops at the first request
that does not). Memory is then a station as well. If every turn needs about $c$ bytes,
the pool holds $m=\lfloor M/c\rfloor$ \emph{slots}; a turn holds its slot from the start
of its prefill to the end of its decode, so the station is a multi-server queue
$\cdot$/G/$m$ whose service time is the holding time $S+D$, and a turn that does not fit
blocks everyone behind it (head-of-line blocking, now on memory). Two caveats: with
unequal contexts a turn needs several slots and the station has no closed form; and in
such engines the cache of paused sessions does not take slots away. Cached blocks that
no running turn uses count as free: an admission that needs them evicts the least
recently used ones, so the slots compete with the cache by turning hits into misses,
not by blocking admission. (The paper's model, which lets a controller keep or drop a
paused state at a price, is a policy these engines do not implement: their KV policy
is LRU on the free blocks.) On the paper's testbed with long contexts this station,
not the prefill queue, set the waits (\cref{L5:sec:honest}). The M/M/$m$ queue is an
exercise of \cref{sec:L2}.

\paragraph{The control changes the service law.} In classical queueing theory the
service-time law is given. Here the KV policy chooses it: a turn whose prefix is still
resident is a hit with short prefill work, and one whose prefix was dropped is a miss
that recomputes the whole context. The policy therefore sets the distribution of $S$ at
the prefill station, and the central question of \cref{sec:L3} is how a performance
measure changes when that distribution changes.

\paragraph{Performance.} The time to first token is the sojourn time at the prefill
station, $\mathrm{TTFT}=W_q+S$. The objective is the mean number of turns in the two
stations, $L=L_P+L_D$; by Little's law (\cref{thm:little}) it is the same as the mean
delay per unit time.

\paragraph{What the projection drops.} Work is measured in seconds rather than tokens
and memory blocks; the budget left to prefill by the decode batch fluctuates, and the
model replaces it by an average; the paper's model has no memory-admission station, so
a turn whose KV does not fit waits for free blocks without the model seeing it; and
session feedback makes turn arrivals non-Poisson.
\Cref{sec:L5} returns to these gaps once the tools are in place.

\section{The exponential distribution and memorylessness}

\begin{definition}[Exponential distribution]\label{L1:def:exp}
A random variable $X\ge0$ with $\Prob(X>t)=e^{-\mu t}$ for $t\ge0$ has the
\term{exponential distribution} with \term{rate} $\mu>0$, written
$X\sim\mathrm{Exp}(\mu)$. Then $\E[X]=1/\mu$ and $\E[X^2]=2/\mu^2$.
\end{definition}

\begin{proposition}[Memorylessness \cite{ross1996}]\label{prop:memoryless}
If $X\sim\mathrm{Exp}(\mu)$, then for all $s,t\ge0$
\begin{equation}\label{L1:eq:memoryless}
  \Prob(X>s+t\mid X>s)=\Prob(X>t).
\end{equation}
Conversely, if a random variable with $0<X<\infty$ a.s.\ satisfies
\eqref{L1:eq:memoryless}, then $X\sim\mathrm{Exp}(\mu)$ for some $\mu\in(0,\infty)$.
\end{proposition}

\begin{proof}
The forward direction is $e^{-\mu(s+t)}/e^{-\mu s}=e^{-\mu t}$.
Converse: let $G(t)=\Prob(X>t)$. Then \eqref{L1:eq:memoryless} reads
$G(s+t)=G(s)G(t)$. $G$ is nonincreasing and right-continuous, and
$G(0)=\Prob(X>0)=1$. By induction $G(nt)=G(t)^n$, hence $G(1)=G(1/n)^n$ and
$G(m/n)=G(1)^{m/n}$.
If $G(1)=0$, then $G(1/n)=0$ for every $n$, which contradicts
$G(1/n)\to G(0)=1$ by right-continuity. If $G(1)=1$, then $G(r)=1$ for every
rational $r$, and by monotonicity $G\equiv1$, that is $X=\infty$ a.s., contradicting
the assumption. Hence $G(1)=e^{-\mu}$ for some $\mu\in(0,\infty)$, and
$G(r)=e^{-\mu r}$ on the rationals. For arbitrary $t$, take rationals $r\downarrow t$;
right-continuity gives $G(t)=e^{-\mu t}$.
\end{proof}

Memorylessness says that the remaining service time of a customer who has already
received $s$ units of service has the same distribution as the service time of a new
customer. This is why, in a system with exponential service, the number in system
$N(t)$ alone determines the future, and $N(t)$ is a continuous-time Markov chain
(\cref{sec:L2}). The following fact is also used often.

\begin{lemma}[Minimum of exponentials \cite{ross1996}]\label{L1:lem:min}
If $X_i\sim\mathrm{Exp}(\mu_i)$, $i=1,\dots,n$, are independent, then
$\min_iX_i\sim\mathrm{Exp}(\sum_i\mu_i)$.
\end{lemma}

\begin{proof}
$\Prob(\min_iX_i>t)=\prod_i\Prob(X_i>t)=e^{-(\sum_i\mu_i)t}$.
\end{proof}

\section{The Poisson process}

\begin{definition}[Counting process and Poisson process]\label{def:poisson}
A \term{counting process} $A=(A(t))_{t\ge0}$ is a nondecreasing, right-continuous,
$\N$-valued process with $A(0)=0$ and jumps of size one only.
A counting process $A$ is a \term{Poisson process} with \term{rate} $\lambda>0$ if
\begin{enumerate}
\item (\term{independent increments}) for disjoint intervals
  $(s_1,t_1],\dots,(s_n,t_n]$ the increments $A(t_j)-A(s_j)$ are independent;
\item (\term{stationary Poisson increments}) for all $s\ge0$, $t>0$,
  $A(s+t)-A(s)\sim\mathrm{Poisson}(\lambda t)$, that is
  $\Prob(A(s+t)-A(s)=n)=e^{-\lambda t}(\lambda t)^n/n!$.
\end{enumerate}
$S_n=\inf\{t:A(t)\ge n\}$ is the $n$-th arrival time, and $T_n=S_n-S_{n-1}$
($S_0=0$) is the $n$-th \term{interarrival time}.
\end{definition}

\begin{theorem}[Exponential interarrival times \cite[\S2.2]{ross1996}]\label{L1:thm:interarrival}
$A$ is a Poisson process with rate $\lambda$ if and only if the interarrival times
$T_1,T_2,\dots$ are independent and each is $\mathrm{Exp}(\lambda)$.
\end{theorem}

\begin{proof}[Proof ($\Rightarrow$)]
For $0=t_0<t_1<\dots<t_n$ and small $h>0$ consider the event
$E_h=\{S_j\in(t_j,t_j+h],\ j=1,\dots,n\}$. Let $F_h$ be the event that there is no
arrival in $(t_{j-1}+h,t_j]$ and exactly one in each $(t_j,t_j+h]$. Then
$F_h\subset E_h$, and $E_h\setminus F_h$ is contained in the event that some interval
of length $h$ has two or more arrivals, which has probability $O(h^{n+1})$. By
independent and Poisson increments,
\[
  \Prob(F_h)=e^{-\lambda(t_1+(t_2-t_1-h)+\cdots+(t_n-t_{n-1}-h))}\,(\lambda h e^{-\lambda h})^n
  =\lambda^n h^n e^{-\lambda t_n}\,(1+O(h)).
\]
Hence $(S_1,\dots,S_n)$ has density $\lambda^ne^{-\lambda s_n}$ on
$\{0<s_1<\dots<s_n\}$. The linear map $T_j=S_j-S_{j-1}$ has Jacobian one and
$s_n=\sum_jT_j$, so $(T_1,\dots,T_n)$ has density $\prod_{j=1}^n\lambda e^{-\lambda t_j}$,
that is, independent $\mathrm{Exp}(\lambda)$. The converse is a standard argument
using that $S_n$ is Gamma distributed together with memorylessness, and we omit it
\cite[\S2.2]{ross1996}.
\end{proof}

\begin{proposition}[Superposition and thinning \cite{kingman1993}]\label{prop:superpose}
\begin{enumerate}[label=(\roman*)]
\item (\term{Superposition}) If $A_1,A_2$ are independent Poisson processes with rates
  $\lambda_1,\lambda_2$, then $A_1+A_2$ is a Poisson process with rate
  $\lambda_1+\lambda_2$.
\item (\term{Thinning}) If each point of a Poisson process with rate $\lambda$ is kept
  with probability $p$ and discarded with probability $1-p$, independently of
  everything else, then the kept points and the discarded points are Poisson
  processes with rates $p\lambda$ and $(1-p)\lambda$, and they are independent.
\end{enumerate}
\end{proposition}

\begin{proof}
(i) Independent increments follow from the independent increments of the two
processes and their mutual independence. For the increment over one interval, the
generating function of independent Poisson variables is
$\E[z^{X_1+X_2}]=e^{\lambda_1t(z-1)}e^{\lambda_2t(z-1)}=e^{(\lambda_1+\lambda_2)t(z-1)}$,
so the sum is $\mathrm{Poisson}((\lambda_1+\lambda_2)t)$.

(ii) Let $X$ be the original increment over an interval of length $t$, $X_1$ the number
kept, and $X_2=X-X_1$ the number discarded. Given $X=j+k$, the conditional
probability that $X_1=j$ is binomial, $\binom{j+k}{j}p^j(1-p)^k$, so
\[
  \Prob(X_1=j,X_2=k)=e^{-\lambda t}\frac{(\lambda t)^{j+k}}{(j+k)!}\binom{j+k}{j}p^j(1-p)^k
  =e^{-p\lambda t}\frac{(p\lambda t)^j}{j!}\cdot e^{-(1-p)\lambda t}\frac{((1-p)\lambda t)^k}{k!}.
\]
Thus on one interval $X_1,X_2$ are independent Poisson variables. On disjoint
intervals the original increments and the coin flips are all independent, so all
increments of the two processes are independent.
\end{proof}

\serving{If sessions arrive as a Poisson process (rate $\Lambda$ in the paper) and each
turn is assigned independently to one of $J$ replicas, then by
\cref{prop:superpose}(ii) each replica also sees Poisson arrivals. The turns within a
session, however, go around a closed loop (prefill $\to$ decode $\to$ tool), so the
turn arrivals are not exactly Poisson. The closed models of \cref{sec:L4} are the tool
for this gap.}

\section{PASTA: Poisson arrivals see time averages}

The distribution of the state seen by arriving customers generally differs from the
time-average distribution. For example, in a D/D/1 system with interarrival times
exactly 1 and service times exactly $0.5$, the server is busy half of the time, yet
every arriving customer finds it idle. With Poisson arrivals this discrepancy
disappears.

\begin{definition}[Lack of anticipation]\label{L1:def:lookahead}
Given a state process $X=(X(t))_{t\ge0}$ and an arrival process $A$, we say that $A$
has \term{lack of anticipation} with respect to $X$ if for every $t$ the future
increments $\{A(t+s)-A(t):s\ge0\}$ are independent of the past
$\{X(u):u\le t\}\cup\{A(u):u\le t\}$.
\end{definition}

In a queueing system this assumption says that future arrivals are independent of the
system state so far. Note the direction: the state may (and does) depend on past
arrivals; what is excluded is that future arrivals react to the state. Poisson arrivals
from outside satisfy it naturally. Arrivals in a closed system (\cref{sec:L4}) do not.

\begin{example}[One process with lack of anticipation, one without]\label{L1:ex:lookahead}
\emph{(a) Lack of anticipation: the M/M/1 queue.} Customers come from a large outside
population as a Poisson process, independently of the service times. The queue length
$X(t)$ is a function of the arrivals before $t$ and the service times, and by
independent increments the arrivals after $t$ are independent of both. An arriving
customer is like a raindrop falling on the queue: the rain does not know how long the
queue is. So the arrival instants are ``random inspection times'', and by
\cref{thm:pasta} the fraction of arrivals that find the server busy equals the fraction
of time it is busy, $\rho$.

\emph{(b) No lack of anticipation: a single agent session.} One session alternates
between a tool call (thinking, mean $Z$) and a turn at the server (mean $1/\mu$). Its
next turn can arrive only after its previous turn has left the server. If $X(t)=1$
(the turn is being served), then no arrival can happen until the service ends; if
$X(t)=0$, an arrival may come at any moment. Future arrivals depend on the current
state, so the condition fails. The consequence is visible: every arrival finds the
server idle, while the server is busy a fraction $(1/\mu)/(Z+1/\mu)>0$ of the time.
Arrivals do not see time averages. With $N$ sessions the same effect is weaker but
present, and \cref{sec:L4} replaces PASTA by the arrival theorem for exactly this
reason.

\emph{(c) Lack of anticipation is not enough.} In the D/D/1 example above, the arrivals
at times $1,2,3,\dots$ are deterministic, hence independent of everything, so they have
lack of anticipation. Yet every arrival finds the server idle. What fails is the other
hypothesis of \cref{thm:pasta}: the arrivals are not Poisson. They are locked in phase
with the service, always landing after the previous customer has left. Poisson
arrivals are spread uniformly over time and cannot lock onto any phase.
\end{example}

\begin{remark}[Why study PASTA if sessions violate it?]\label{L1:rem:why-pasta}
Example~\ref{L1:ex:lookahead}(b) shows that a single agent session does not have lack
of anticipation. PASTA is still the right starting point, for three reasons.
\begin{enumerate}
\item \emph{The prefill model rests on it.} The paper models the prefill stage as an
  M/G/1 queue, and the proof of the Pollaczek--Khinchine formula (\cref{thm:pk}) uses
  PASTA. The paper states that session feedback makes this an approximation and tests
  it. One cannot judge an approximation without knowing what it approximates.
\item \emph{Many sessions look Poisson.} Each session reacts to the state, but the
  superposition of many independent sessions, each contributing a small share of the
  arrivals, is close to a Poisson process; this is the Palm--Khinchine theorem
  \cite{khinchine1960}. New sessions themselves arrive as a Poisson process in the
  model, so the M/G/$\infty$ and Campbell calculations of \cref{sec:L4} use Poisson
  arrivals without approximation.
\item \emph{The exact closed-system answer is a corrected PASTA.} The arrival theorem
  (\cref{thm:arrival}) says that an arriving turn sees the time averages of the system
  with one session fewer, and it reduces to PASTA as the number of sessions grows
  (\cref{L4:rem:arrival-limit}). For exponential work, \cref{prop:finite}(ii) shows
  that the open formula built on PASTA overstates the wait at equal utilisation, so
  PASTA gives a conservative estimate whose error the finite-population model
  quantifies.
\end{enumerate}
\end{remark}

\begin{theorem}[PASTA \cite{wolff1982}]\label{thm:pasta}
Let $A$ be a Poisson process with rate $\lambda$, $X$ a state process with left
limits, and $B$ a set of states, and suppose $A$ has lack of anticipation with respect
to $X$. If $U(t)=\1\{X(t^-)\in B\}$ is left-continuous, then for every $T>0$
\begin{equation}\label{L1:eq:pasta}
  \E\Bigl[\sum_{k:\,S_k\le T}\1\{X(S_k^-)\in B\}\Bigr]=\lambda\,\E\Bigl[\int_0^T\1\{X(t)\in B\}\dd t\Bigr].
\end{equation}
Consequently, if the limits of the two averages
$\frac1T\int_0^T\1\{X(t)\in B\}\dd t$ and
$\frac{1}{A(T)}\sum_{k\le A(T)}\1\{X(S_k^-)\in B\}$ exist with probability one, they
are equal: the fraction of arrivals that find the system in $B$ equals the fraction of
time it spends in $B$.
\end{theorem}

\begin{proof}
The left side is $\int_0^TU(t)\dd A(t)$. For the partition $t_j=jT/n$ of $[0,T]$,
define the step function $U_n(t)=\sum_{j}U(t_j)\1\{t\in(t_j,t_{j+1}]\}$. $U(t_j)$ is
determined by the past up to time $t_j$, and by lack of anticipation the increment
$A(t_{j+1})-A(t_j)$ is independent of that past, so
\[
  \E\Bigl[\int_0^TU_n\dd A\Bigr]=\sum_j\E[U(t_j)]\,\E[A(t_{j+1})-A(t_j)]
  =\lambda\sum_j\E[U(t_j)]\,(t_{j+1}-t_j)=\lambda\,\E\Bigl[\int_0^TU_n(t)\dd t\Bigr].
\]
Since $U$ is left-continuous, $U_n(t)\to U(t)$ for every $t$. The integrands on the
two sides are bounded by $A(T)$ and $T$ respectively, and $\E[A(T)]=\lambda T<\infty$,
so by dominated convergence $\E[\int_0^TU\dd A]=\lambda\E[\int_0^TU\dd t]$ as
$n\to\infty$. Finally, the set of $t$ with $X(t^-)\ne X(t)$ is countable, so replacing
$U(t)$ by $\1\{X(t)\in B\}$ on the right does not change the integral.
The statement about limits follows from a sample-path version of
\eqref{L1:eq:pasta} (a law of large numbers applied to a martingale), for which we
cite \cite{wolff1982}.
\end{proof}

In \cref{sec:L3}, PASTA becomes a key step of the Pollaczek--Khinchine formula in the
form ``the mean work an arriving turn finds waiting equals the time-average work
waiting''.

\section{The renewal-reward theorem}

\begin{definition}[Renewal process and rewards]\label{L1:def:renewal}
For i.i.d.\ positive random variables $C_1,C_2,\dots$ with $\E[C_1]<\infty$, let
$\tau_n=C_1+\dots+C_n$ and $M(t)=\#\{n\ge1:\tau_n\le t\}$. $M$ is a
\term{renewal process} and $C_n$ is its $n$-th \term{cycle}. Suppose each cycle
carries a reward $R_n\ge0$, the pairs $(C_n,R_n)$ are i.i.d., and the cumulative reward
$R(t)$ up to time $t$ satisfies
\[
  \sum_{n=1}^{M(t)}R_n\;\le\;R(t)\;\le\;\sum_{n=1}^{M(t)+1}R_n
\]
(this includes rewards that accrue gradually within a cycle). This is a
\term{renewal-reward process}.
\end{definition}

\begin{theorem}[Renewal-reward theorem \cite{ross1996}]\label{thm:renewal-reward}
If $\E[R_1]<\infty$, then with probability one
$\displaystyle\lim_{t\to\infty}\frac{R(t)}{t}=\frac{\E[R_1]}{\E[C_1]}$.
\end{theorem}

\begin{proof}
First we show $M(t)/t\to1/\E[C_1]$ a.s. By definition $\tau_{M(t)}\le t<\tau_{M(t)+1}$,
and since each $\tau_n<\infty$, $M(t)\to\infty$. By the strong law,
$\tau_n/n\to\E[C_1]$, so both ends of
\[
  \frac{\tau_{M(t)}}{M(t)}\le\frac{t}{M(t)}<\frac{\tau_{M(t)+1}}{M(t)+1}\cdot\frac{M(t)+1}{M(t)}
\]
converge to $\E[C_1]$. Now
\[
  \frac{1}{t}\sum_{n=1}^{M(t)}R_n=\frac{M(t)}{t}\cdot\frac{1}{M(t)}\sum_{n=1}^{M(t)}R_n
  \longrightarrow\frac{1}{\E[C_1]}\,\E[R_1]
\]
(the strong law for the second factor), and the upper bound has the same limit since
$(M(t)+1)/t\to1/\E[C_1]$. The conclusion follows by squeezing.
\end{proof}

In the form ``long-run time average $=$ expected reward per cycle / expected cycle
length'', the renewal-reward theorem is used directly for the residual service time
(the inspection paradox) in \cref{sec:L3}.

\section{Little's law}

\begin{theorem}[Little's law \cite{little1961,stidham1974}]\label{thm:little}
For a sequence of customers $(a_k,d_k)$ with $a_k\le d_k<\infty$ and $a_k$
nondecreasing, suppose the two limits
\[
  \lambda=\lim_{t\to\infty}\frac{A(t)}{t}\in(0,\infty),\qquad
  W=\lim_{n\to\infty}\frac1n\sum_{k=1}^nW_k<\infty
\]
exist. Then $L=\lim_{t\to\infty}\frac1t\int_0^tN(s)\dd s$ exists and
\begin{equation}\label{L1:eq:little}
  L=\lambda W .
\end{equation}
\end{theorem}

\begin{proof}
Since $N(s)=\sum_k\1\{a_k\le s<d_k\}$, we have
$\int_0^tN(s)\dd s=\sum_k\bigl|[a_k,d_k)\cap[0,t]\bigr|$.
Each term is at most $W_k$, is zero if $a_k>t$, and equals $W_k$ exactly if
$d_k\le t$. Therefore
\begin{equation}\label{L1:eq:little-sandwich}
  \sum_{k:\,d_k\le t}W_k\;\le\;\int_0^tN(s)\dd s\;\le\;\sum_{k:\,a_k\le t}W_k=\sum_{k=1}^{A(t)}W_k .
\end{equation}
Upper bound: since $A(t)\to\infty$,
$\frac1t\sum_{k\le A(t)}W_k=\frac{A(t)}{t}\cdot\frac{1}{A(t)}\sum_{k\le A(t)}W_k\to\lambda W$.

Lower bound: with $\bar W_n=\frac1n\sum_{k\le n}W_k$ we have
$\frac{W_n}{n}=\bar W_n-\frac{n-1}{n}\bar W_{n-1}\to W-W=0$. Also $A(a_n)\ge n$ and
$A(a_n^-)\le n-1$, so $a_n/n\to1/\lambda$, and therefore $W_n/a_n\to0$.
For any $\varepsilon>0$ there is $k_0$ such that $d_k=a_k+W_k\le(1+\varepsilon)a_k$ for
$k\ge k_0$. Then for $t\ge\max_{k<k_0}d_k$, every customer with
$a_k\le t/(1+\varepsilon)$ has $d_k\le t$, so
\[
  \frac1t\int_0^tN(s)\dd s\;\ge\;\frac1t\sum_{k=1}^{A(t/(1+\varepsilon))}W_k
  \longrightarrow\frac{\lambda W}{1+\varepsilon}.
\]
Letting $\varepsilon\downarrow0$, the liminf is at least $\lambda W$. Together with
the upper bound this gives \eqref{L1:eq:little}.
\end{proof}

The proof uses no distributional assumption and no service discipline. It is a
deterministic theorem that holds on every sample path on which the limits exist. So
depending on how we draw the boundary of ``the system'', we obtain several identities.

\begin{corollary}[Utilisation law \cite{kleinrock1975}]\label{cor:busy}
Under the assumptions of \cref{thm:little}, consider a single-server system and
suppose the average of the service times $\frac1n\sum_{k\le n}S_k$ converges to
$\E[S]$.
\begin{enumerate}[label=(\roman*)]
\item The fraction of time the server is busy is $\rho=\lambda\E[S]$. In particular
  $\rho\le1$ in a stable system.
\item For the waiting room, $L_q=\lambda W_q$, and under FIFO $L=L_q+\rho$.
\end{enumerate}
\end{corollary}

\begin{proof}
(i) Take ``the server'' alone as the system. Customer $k$ enters it when its service
starts and leaves $S_k$ later. Every customer is eventually served, so the arrival
rate of this subsystem is also $\lambda$, and the customer average of its sojourn times
is $\E[S]$. The number in this subsystem is 1 when the server is busy and 0 otherwise,
so its time average is the fraction of time the server is busy. By \cref{thm:little}
this equals $\lambda\E[S]=\rho$, and being a fraction it is at most one.
(ii) Taking the waiting room alone as the system, the sojourn times are $W_{q,k}$, so
$L_q=\lambda W_q$. Taking time averages in $N=N_q+(\text{number in service})$ gives
$L=L_q+\rho$.
\end{proof}

\begin{example}\label{L1:ex:little}
Suppose $\lambda=2$ turns per second arrive at the prefill stage and the mean prefill
work is $\E[S]=0.3$ seconds. By \cref{cor:busy} the server is busy a fraction
$\rho=0.6$ of the time. If the observed mean number of waiting turns is $L_q=0.9$,
then $W_q=L_q/\lambda=0.45$ seconds, the mean TTFT is $W_q+\E[S]=0.75$ seconds, and the
mean number of turns in the stage is $L=\lambda\cdot0.75=1.5$.
\end{example}

The paper takes as its objective the mean number of turns in the prefill stage and in
the decode batch, $L=L_P+L_D$. By \cref{thm:little} this equals the total delay
accumulated per unit time, $\lambda\cdot(\text{mean sojourn time})$, so ``reduce the
mean number of turns'' and ``reduce the mean delay'' are the same goal. The ``price of
a miss'' of \cref{sec:L3} is defined precisely as a change in this $L_P$.

\section{Exercises}

\begin{exercise}\label{L1:exr:race}
For independent $X_1\sim\mathrm{Exp}(\mu_1)$ and $X_2\sim\mathrm{Exp}(\mu_2)$, show that
$\Prob(X_1<X_2)=\mu_1/(\mu_1+\mu_2)$, and that the event $\{X_1<X_2\}$ is independent
of $\min(X_1,X_2)$.
\emph{Hint:} integrate $\Prob(X_1<X_2,\ \min>t)$ against the density of $X_1$.
\end{exercise}

\begin{exercise}\label{L1:exr:delay}
Each time a session makes a tool call, it stays away from the server for a mean of $Z$
seconds. If tool calls start at rate $\lambda$, show that the time-average number of
sessions in a tool call is $\lambda Z$. Does this hold whatever the distribution of the
tool-call duration?
\emph{Hint:} take the tool-call stage alone as ``the system'' and apply
\cref{thm:little}.
\end{exercise}
